% Formalization support lemmas.
%
% NOT book material. These declarations were invented in this workspace to
% carry a Lean formalization and have no counterpart in the source text --
% they carry no \dcref. They live here, outside the book chapters, so the
% book blueprint stays timeless mathematics while their \leanok marks and
% Lean anchors are preserved.
%
% Do not add book mathematics to this file, and do not add formalization
% notes to the book chapters.

\begin{definition}[Isometry]
\label{def:isometry}
\uses{def:local-isometry}
\lean{MorganTianLib.LocalIsometry.IsIsometry}
\leanok
An {\sl isometry} is a bijective local isometry.
\end{definition}

\begin{definition}[Distributional Laplacian inequality]
\label{def:weak-laplacian}
\uses{def:riemannian-metric}
\lean{MorganTianLib.WeakLaplacianLE}
\leanok
Fix an additive Haar normalization for the volume measure of a Riemannian
manifold $(M,g)$.  For a locally Lipschitz function $f$ and a locally
integrable function $h$, write $\triangle f\le h$ in the distributional
(weak) sense when, for every non-negative compactly supported smooth test
function $\phi$,
$$\int_M f\triangle \phi\,d\vol_g\le
\int_M h\phi\,d\vol_g.$$
\end{definition}

\begin{lemma}[Abstract cone curvature block]
\label{lem:abstract-cone-curvature-block}
\uses{def:curvature-operator}
\lean{MorganTianLib.coneCurvatureBlock,
  MorganTianLib.coneCurvatureBlock_apply,
  MorganTianLib.coneCurvatureBlock_eigenpair}
\leanok
Let $A$ be a bounded endomorphism of a normed real vector space $V$ and define
$B_s$ on $V\oplus W$ by
$$B_s(x,y)=(s^2Ax,0).$$
If $Ax=\lambda x$, then
$B_s(x,0)=s^2\lambda(x,0)$, while $B_s(0,y)=0$ for every $y\in W$.
\end{lemma}

\begin{lemma}[Nonsingularity of the constant-curvature exponential map]
\label{lem:constant-curvature-exp-nonsingular-nonzero}
\uses{lem:constant-curvature-normal-jacobi-profile,
  def:conjugate-point,def:exponential-map,thm:hopf-rinow}
\lean{MorganTianLib.constantCurvature_expMapGlobal_localDiffeomorph}
\leanok
Let $(M,g)$ be complete of constant sectional curvature $K$.
If $v\in T_pM$ is nonzero and $K|v|^2<\pi^2$, then $\exp_p$ is a
local diffeomorphism at $v$.
\end{lemma}

\begin{lemma}[Constant-curvature normal Jacobi profile]
\label{lem:constant-curvature-normal-jacobi-profile}
\uses{lem:constant-curvature-tensor,def:geodesic,lem:parallel-frame}
\lean{MorganTianLib.exists_constantCurvatureJacobiProfile,
  MorganTianLib.constantCurvatureJacobi_normal_profile}
\leanok
For every $K\in\mathbb R$ there are functions $s,s'$ satisfying
$s''+Ks=0$, $s(0)=0$, and $s'(0)=1$.  Along a unit-speed geodesic
in constant curvature $K$, if $E$ is parallel and perpendicular to
the velocity, then $J=sE$ is a Jacobi field with covariant derivative
$\nabla J=s'E$.
\end{lemma}

\begin{lemma}[Connection and curvature under constant rescaling]
\label{lem:constant-rescaling-connection-curvature}
\uses{def:riemannian-metric,thm:levi-civita-connection,
  def:riemann-curvature-tensor,def:sectional-curvature}
\lean{MorganTianLib.rescaledMetric,
  MorganTianLib.rescaledMetric_metricInner,
  MorganTianLib.AffineConnection.IsMetricCompatible.rescaledMetric,
  MorganTianLib.rescaledMetric_leviCivitaConnection,
  MorganTianLib.rescaledMetric_curvature,
  MorganTianLib.rescaledMetric_curvatureForm,
  MorganTianLib.rescaledMetric_isConstantCurvature}
\leanok
For $c>0$, the metric $g_c=cg$ has pairing
$\langle v,w\rangle_{g_c}=c\langle v,w\rangle_g$.  Its Levi-Civita
connection and $(1,3)$ curvature tensor agree with those of $g$, while
its $(0,4)$ curvature form is multiplied by $c$.  If $g$ has constant
sectional curvature $K$, then $g_c$ has constant sectional curvature
$K/c$.
\end{lemma}

\begin{lemma}[Analytic reduction for the distance-function bound]
\label{lem:distance-function-weak-laplacian-reduction}
\uses{def:weak-laplacian,rem:laplacian-weak-green-identity}
\lean{MorganTianLib.DistanceWeakGreenIdentity,
  MorganTianLib.DistancePolarTestInequality,
  MorganTianLib.lipschitzWith_one_distanceFrom,
  MorganTianLib.locallyLipschitz_distanceFrom,
  MorganTianLib.weakLaplacianLE_distanceFrom_of_green_and_polar}
\leanok
The function $x\mapsto d(p,x)$ is $1$-Lipschitz, hence locally
Lipschitz.  Suppose $(n-1)/d(p,\cdot)$ is locally integrable and
that a chosen radial pairing satisfies both the weak Green identity
and the polar-coordinate test inequality.  Then
$$\triangle d(p,\cdot)\le \frac{n-1}{d(p,\cdot)}$$
in the weak sense.
\end{lemma}

\begin{lemma}[Induced connection on an identified immersed patch]
\label{lem:immersed-patch-induced-connection}
\uses{def:riemannian-metric}
\lean{MorganTianLib.submanifoldConnection_eq_tangentProjection}
\leanok
For an identified immersed patch $D$, an ambient affine connection
$\nabla$, and vector fields $X,Y$, the induced covariant derivative is
the tangential projection of the ambient one:
$$D.\operatorname{inducedCov}_{\nabla}(X,Y)
=D.\operatorname{tangentProj}(\nabla_XY).$$
\end{lemma}

\begin{lemma}[Complete connected-source local isometries are coverings]
\label{lem:local-isometry-covering-connected-source}
\uses{def:local-isometry,lem:local-isometry-exponential,thm:hopf-rinow}
\lean{MorganTianLib.LocalIsometry.IsLocalIsometry.surjectiveCovering}
\leanok
Let $F\colon N\to M$ be a local isometry between connected Riemannian
manifolds.  If $N$ is complete, then $F$ is a surjective covering map.
\end{lemma}

\begin{lemma}[Matrix-Jacobi Rauch lower bound]
\label{lem:matrix-jacobi-rauch-lower}
\uses{lem:vector-sturm-comparison,lem:radial-shape-riccati}
\lean{MorganTianLib.rauch_lower_matrix_norm,
  MorganTianLib.rauch_lower_matrix_sq}
\leanok
Let $\mathcal J$ be radial matrix-Jacobi data on $[0,b]$.  Suppose
$K\ge0$, $T\le b$, $\sqrt K\,T<\pi$, and
$\langle\mathcal R(t)X,X\rangle\le K|X|^2$ for $0<t<T$.
Then, for $0<t\le T$ and every $w$,
$$|w|\,s_K(t)\le|\mathcal J(t)w|,
\qquad s_K(t)^2|w|^2\le|\mathcal J(t)w|^2.$$
\end{lemma}

\begin{lemma}[Radial structure of the metric segment domain]
\label{lem:metric-segment-domain-radial}
\uses{def:cut-time-locus,prop:minimal-geodesic-no-conjugate}
\lean{MorganTianLib.cutTime_radial_interval,
  MorganTianLib.isMinimizingUpTo_one_of_mem_segmentDomain,
  MorganTianLib.eq_of_mem_segmentDomain_of_isMinimizingUpTo_one,
  MorganTianLib.dist_expMapGlobal_eq_norm_of_mem_segmentDomain,
  MorganTianLib.radial_minimizing_and_no_conjugate_before_cut,
  MorganTianLib.cutLocus_properties_metric_core}
\leanok
Let $c(v)$ be the metric cut time and let
$U_p^{\mathrm{met}}=\{v\in T_pM\mid c(v)>1\}$.  For every $t>0$,
$$tv\in U_p^{\mathrm{met}}\quad\Longleftrightarrow\quad t<c(v).$$
Every $w\in U_p^{\mathrm{met}}$ gives the unique minimizing radial geodesic
from $p$ to $\exp_p(w)$ and $d(p,\exp_p(w))=|w|$.  A unit radial geodesic is
minimizing and has no conjugate point at every positive time strictly before
its cut time.  With a Haar normalization on the tangent space, the
corresponding metric cut locus has zero Riemannian measure.
\end{lemma}

\begin{lemma}[Pointwise normal and strongly convex balls]
\label{lem:pointwise-strongly-convex-normal-ball}
\uses{def:exponential-map,def:geodesic}
\lean{MorganTianLib.exists_normal_ball_core,
  MorganTianLib.exists_stronglyConvex_normal_ball}
\leanok
For every $p\in M$ there is a positive radius on whose model-space
ball the exponential map is defined and injective, with open image.
There is also $\beta>0$ such that the closed metric ball
$\overline B(p,\beta)$ is strongly convex: every two of its points
are joined by a unique minimizing geodesic whose open arc stays in
the ball.
\end{lemma}

\begin{remark}[Cut-time injectivity radius as distance to the metric cut locus]
\label{rem:cut-time-injectivity-radius-distance}
\uses{def:cut-time-locus}
\lean{MorganTianLib.cutLocusDistance,
  MorganTianLib.injectivityRadius_eq_cutLocusDistance}
\leanok
The infimum of the metric cut times over the unit tangent sphere equals the
extended distance from $p$ to the corresponding metric cut locus; both are
$+\infty$ when that cut locus is empty.
\end{remark}

\begin{remark}[The end compactification]
\label{rem:ends-compactification}
\uses{def:space-of-ends}
The space of ends of a connected manifold is a compact space, and
one can define a topology on the union of $M$ and its space of ends
making this union a compact, connected Hausdorff space which is a
compactification of $M$. We will not use these facts in this
chapter.
\end{remark}

\begin{remark}[Injectivity radius as distance to the cut locus]
\label{rem:injectivity-radius-distance}
\uses{def:injectivity-radius,def:cut-locus}
\lean{MorganTianLib.cutLocusDistance,
  MorganTianLib.injectivityRadius_eq_cutLocusDistance}
\leanok
The injectivity radius $\inj_M(p)$ is the distance in $M$ from $p$ to
the cut locus ${\mathcal C}_p$, with the value $+\infty$ when the cut locus is empty.
\end{remark}

\begin{remark}[Injectivity radius as distance to the segment-domain frontier]
\label{rem:injectivity-radius-frontier}
\uses{def:injectivity-radius,def:cut-locus}
\lean{MorganTianLib.injectivityRadius_eq_segmentDomainFrontierDistance}
\leanok
The injectivity radius $\inj_M(p)$ is the
distance in $T_pM$ from $0$ to the frontier of $U_p$.
\end{remark}

\begin{remark}[Rademacher and the weak Green identity]
\label{rem:laplacian-weak-green-identity}
\uses{lem:laplacian-local-formula,rem:laplacian-weak-sense}
By Rademacher's theorem a locally Lipschitz function is
differentiable almost everywhere, its almost-everywhere defined
gradient $\nabla f$ is locally bounded, and it coincides with the
distributional gradient; in particular the restriction of $\nabla
f$ to any compact subset of $M$ is an $L^2$ one-form. Moreover, for
every test function $\phi$,
$$\int_Mf\triangle
\phi\, d\vol=-\int_M\langle \nabla f,\nabla\phi\rangle\, d\vol.$$
Indeed, covering the compact support of $\phi$ by finitely many
charts and using a subordinate partition of unity, it suffices to
prove this in a single chart; there, by the divergence form of the
Laplacian (Lemma~\ref{lem:laplacian-local-formula}),
$$\sqrt{\det g}\ \triangle\phi
=\partial_i\left(\sqrt{\det g}\,g^{ij}\,\partial_j\phi\right),$$
and the assertion becomes the integration-by-parts identity for the
locally Lipschitz function $f$ against the compactly supported
smooth vector field with components $\sqrt{\det
g}\,g^{ij}\partial_j\phi$, which is the statement that the
almost-everywhere gradient of $f$ represents its distributional
gradient (standard real analysis).
\end{remark}

\begin{remark}[Local inverses of local isometries]
\label{rem:local-isometry-local-inverse}
\uses{def:local-isometry,def:isometry}
\lean{MorganTianLib.LocalIsometry.IsLocalIsometry.exists_isometry_restriction,
  MorganTianLib.LocalIsometry.IsIsometry.symm}
\leanok
By the inverse function theorem, a local isometry restricts to an
isometry from a sufficiently small neighborhood of each point onto
its image.  The inverse of an isometry is again an isometry.
\end{remark}

\begin{remark}[Agreement of the metric and book cut-locus definitions]
\label{rem:metric-book-cut-locus-agreement}
\uses{def:cut-time-locus,def:cut-locus,def:injectivity-radius}
% No \leanok: the original bookInjectivityRadius uses the fixed model norm,
% while injectivityRadius uses the metric norm g.metricInner; the corrected
% metric-ball object and conditional norm-compatibility bridge are recorded
% in MorganTianLib.InjectivityRadiusAgreement.
The metric objects of Definition~\ref{def:cut-time-locus} agree with
the $U_p$, ${\mathcal C}_p$ and $\inj_M(p)$ of
Definitions~\ref{def:cut-locus} and~\ref{def:injectivity-radius}.
\end{remark}

\begin{remark}[Provenance of the regularity theorem]
\label{rem:weak-harmonic-regularity-citation}
\uses{thm:weak-harmonic-regularity,claim:dirichlet-problem-small-balls,claim:weak-subsolution-comparison,rem:laplacian-weak-sense,lem:laplacian-local-formula}
Theorem~\ref{thm:weak-harmonic-regularity} is Weyl's lemma for the
Laplace--Beltrami operator: in local coordinates $\Delta$ is a
second-order elliptic operator with smooth coefficients
(Lemma~\ref{lem:laplacian-local-formula}), and continuous
distributional solutions of such equations are smooth. The proof
given above derives it from the two citation-backed claims of this
section, which are therefore the only analytic black boxes that
this chapter consumes. For Riemannian manifolds the theorem is
stated as the regularity theorem for harmonic functions in
\cite{Petersen}, p.~279ff.\ (in the third edition: Theorem 7.1.8,
pp.~283--284), and proved there by harmonic replacement; note,
however, that Petersen's development of the maximum principle is
in the barrier sense, so his argument does not transpose verbatim
to the distributional setting of
Remark~\ref{rem:laplacian-weak-sense}. For the underlying PDE
facts see Gilbarg--Trudinger, {\it Elliptic Partial Differential
Equations of Second Order}, Chapters 6 and 8, and H.~Weyl, {\it
The method of orthogonal projection in potential theory}, Duke
Math.~J.\ {\bf 7} (1940), 411--444.
\end{remark}
